% !TeX root = ../thesis.tex
\chapter{绪论}
\label{chap:ch01_introduction}

\section{研究背景与意义}
\label{sec:research_background}

机器人灵巧操作需要在物体、手指与环境持续交互的过程中实现运动和接触的协调。与将物体从一个位置搬运到另一位置相比，调整物体姿态、稳定持握以及沿约束方向完成装配，对局部接触状态和运动精度提出了不同要求。多指灵巧手能够通过改变接触位置与手指配合方式完成多样化操作，但也需要处理多个接触点之间的相互作用。已有手内操作研究展示了换指、多指协调及利用重力调整物体姿态等行为，体现了这类任务中运动控制与接触变化的紧密联系\cite{openai2020dexterous}。对于持续接触和多指协调任务，预先给定的末端位姿通常还需要与观测和执行反馈相结合，以估计当前状态、选择动作并修正后续运动。如何获得与接触过程相关的感知信息，并使其服务于可执行的操作，是研究灵巧手操作的一项基础问题。

视觉能够提供物体外形、空间位置与周围环境信息，是操作规划的重要依据。当机械臂和手指接近物体时，遮挡可能削弱对局部接触区域的观察；接触已经发生并不意味着物体能够保持稳定，初始位姿相近的抓取也可能因接触条件不同而产生不同结果。触觉通过直接接触获取局部信息，可为接触状态识别、物体属性判断及滑移检测提供线索，与视觉观测形成补充\cite{luo2017tactile}。Calandra等将视觉、触觉与候选动作共同用于预测抓取调整后的结果，并据此选择重抓取动作，说明接触后的反馈具有参与动作调整的价值\cite{calandra2018feeling}。这种价值并不意味着触觉可以独立替代视觉，而是表明不同模态对应着不同的可观测信息。对于多指接触，还需要同时关注响应发生的位置、强弱及其随时间的变化，才能进一步研究哪些感知特征与持握、移动和释放等操作阶段有关。

对上述问题开展研究，需要能够反复执行并比较操作过程的实验条件。真实机器人试验涉及设备使用、物体复位与维护，接触操作还受到碰撞和施力范围等安全约束。多指操作研究中对实际试验耗时与成本的讨论，以及接触丰富示教采集中的安全问题，反映了扩大实机试验规模的现实困难\cite{openai2020dexterous,fang2024rh20t}。仿真能够提供可配置的场景和初态，便于在受控条件下生成轨迹、记录观测并分析失败，但其对真实系统的近似仍需单独检验\cite{tobin2017randomization}。在仿真中，具有较好几何评价的抓取候选也不等同于机器人能够到达并完成的抓取；已有抓取研究将候选获取与可达性筛选区分开来\cite{bohg2014grasp}。本文搭建的仿真平台采用RM75B机械臂与途见科技（北京）有限公司（以下简称途见科技）自研的欠驱动灵巧手模型，需要将物体相对抓取姿态转化为机械臂运动，并检查接近、闭合、保持和后续任务执行。因此，候选生成、机器人可达性筛选与轨迹及任务验证需要相互衔接，才能获得可供使用的专家示教。

专家示教的使用还依赖数据本身的组织与质量。机器人状态描述执行条件，动作记录控制指令，视觉反映场景及物体变化，触觉则补充局部接触信息；这些数据只有与具体执行过程建立明确关联，才能用于分析观测与动作之间的关系。关于离线示教学习的研究表明，示教质量、观测选择与历史信息都会影响策略表现，数据条数增加并不能替代对数据内容的检查\cite{mandlekar2022demonstrations}。BridgeData V2通过组织不同任务、环境和采集来源支持操作学习，RH20T则将多模态记录的时间戳、坐标标定与数据处理纳入采集流程\cite{walke2023bridgedata,fang2024rh20t}。本项目已有抓取抬升、抓取放置、堆叠、推动和螺母装配任务的开发记录，并进行了LeRobot格式的数据构建。对于这些来源和版本不同的数据，需要检查记录发生在动作执行前还是执行后、采样频率是否一致，以及转换后是否保留原始执行索引。同一次仿真执行内部的数据关联，也应与仿真和实机之间的同步采集区别对待，不能仅凭相同的字段或时间戳形式认定两域已经建立物理配对。

即使操作和记录链路能够运行，仿真触觉也不能仅凭输出维度与实机一致，就被视为具有相同含义的观测。传感器结构、接触到响应的转换方式，以及滤波、归一化等处理，都会影响最终信号的表达。TACTO面向视觉触觉传感器，将物理仿真中的接触与传感器成像联系起来；Taxim则针对GelSight的光学响应和标记位移建立模型，并利用实测数据进行校准\cite{wang2022tacto,si2022taxim}。这些工作体现了触觉模型对具体传感机制和读出假设的依赖，不能直接作为本项目阵列观测的物理等价依据。本项目的仿真端通过MuJoCo原生触觉传感器及后续信号处理获得712维标量观测，实机端使用既有处理后触觉记录，目前缺少经过认证的同步物理接触配对。因此，真实与仿真触觉的跨域重建既要研究分布差异，也要关注输入接触结构能否保留。利用真实触觉数据的统计特征构造的伪配对可以提供有限的对应监督，但不能直接作为同一次物理接触的真值；其适用范围和重建效果仍需通过实验评价。

基于这些需求，本文将仿真平台、专家示教、数据质量与触觉跨域重建放在相互关联的研究过程中考察。机器人与触觉观测的集成，为在统一仿真条件下分析多任务接触过程提供实验载体；将抓取候选与机械臂约束、轨迹执行相结合，有助于明确几何抓取方案转化为可执行专家示教所需满足的条件。对多模态数据来源、关联和转换过程进行审计，则使后续分析能够追溯到具体执行记录，减少数据语义不一致带来的歧义。在此基础上，利用真实统计信息探索仿真触觉重建，研究分布逼近与输入结构保持之间的关系，可以帮助认识仿真数据用于真实触觉研究时的条件与局限。这些工作的工程价值在于连接模型、操作与数据流程，方法研究价值则在于分析可执行性、监督对应和重建质量之间的联系，而不在于将平台集成或某一评价指标的改善直接等同于真实操作能力的提升。

\section{国内外研究现状}
\label{sec:related_work}

\subsection{灵巧手抓取生成与可执行操作}
\label{subsec:grasp_generation_execution}

抓取生成需要根据物体几何和手部结构，确定接触位置、手掌位姿及手指配置。解析方法通常以力闭合、接触约束和穿透惩罚等指标评价候选，数据驱动方法则利用已有抓取样本学习从物体观测到抓取配置的关系。Bohg等的综述按照物体先验和数据来源梳理了相关路线，并指出抓取选择还需考虑机器人可达性\cite{bohg2014grasp}。对于多指手，接触组合和关节配置相互影响，生成过程既要避免手物穿透，又要保留不同抓持方式。北京大学等机构的DexGraspNet利用可微力闭合估计与几何能量优化生成抓取，再通过物理仿真检验候选，为学习式抓取提供数据\cite{wang2023dexgraspnet}。这种路线将抓取搜索与数据构建联系起来，但其评价仍与所采用的手型、接触假设和仿真测试有关，不能直接解释为任意机械臂系统上的任务成功率。

优化目标如何近似真实抓持，是候选生成中的关键问题。GraspQP将接触扳手的受限组合写为隐式二次规划，以构建可微的力闭合能量，并结合调整后的采样优化过程兼顾候选质量与多样性\cite{zurbrugg2025graspqp}。与依赖解析能量的搜索不同，DexEvolve在解析候选基础上引入仿真反馈，通过选择、变异和交叉细化抓取，随后将所得分布蒸馏为基于点云的生成模型\cite{zurbrugg2026dexevolve}。两者体现了不同的处理思路：解析优化便于定义搜索方向，物理反馈则能够检验并修正部分简化假设带来的偏差。物理检验也并非没有条件，其结果受到接触参数、扰动方式和执行控制的影响。因此，比较这些方法时，需要区分候选的几何评价、实际抓持稳定性以及候选多样性，不能仅以同名的成功率指标判断优劣。

从稳定抓持进一步走向可执行操作，还需要建立目标姿态与运动过程之间的联系。UniDexGrasp将多样抓取候选生成与目标条件抓取策略分开，通过姿态分布建模、课程学习和教师学生策略学习连接视觉观测与抓取执行\cite{xu2023unidexgrasp}。这表明候选生成与执行能够协同设计，但两者仍有各自的约束。对于机械臂与欠驱动灵巧手组成的系统，手部配置需要符合驱动耦合关系，腕部目标需要满足机械臂运动学约束，接近和搬运过程还需检查碰撞及任务要求。本文的RM75B机械臂与途见科技灵巧手仿真平台据此将开源抓取方法的适配、可达性筛选与执行验证相衔接。第三方方法提供候选基础，本文的适配工作侧重具体平台上的约束连接与专家轨迹构建。其中，目标位姿的可达性检查与完整轨迹验证承担不同职责，需要分别给出依据。

\subsection{机器人仿真与触觉感知建模}
\label{subsec:simulation_tactile_modeling}

机器人仿真为控制和操作研究提供动力学、碰撞检测与接触计算基础。MuJoCo采用广义坐标描述多关节系统，并将接触和约束求解纳入动力学计算，支持执行器及机构约束建模\cite{todorov2012mujoco}。面向大规模策略学习，Isaac Gym通过GPU物理仿真和张量接口连接环境状态与学习过程，减少仿真与训练之间的数据传输开销\cite{makoviychuk2021isaacgym}。这类平台主要解决物体和机器人如何运动、如何发生接触的问题，而特定触觉传感器的输出还涉及接触量如何作用于敏感材料，以及形变或受力如何转换为信号。动力学求解器提供的接触信息并不能自动替代传感器模型；仿真管线的计算效率也不能直接说明触觉观测的物理精度。

触觉模拟因传感机制不同而呈现不同路线。TACTO将物理场景同步与光学渲染相连接，通过传感器的相机、光照和表面配置生成视觉触觉图像，并可利用真实背景图像修正部分外观差异\cite{wang2022tacto}。Taxim针对GelSight建立实例驱动的光学模型，并利用有限元标定的形变关系模拟标记位移\cite{si2022taxim}。前者侧重灵活渲染，后者进一步结合特定传感器的校准与形变传播，两者均依赖各自的几何和读出假设。对曲面阵列而言，还需明确感知单元的位置、朝向与覆盖区域，以及空间采样后信号处理的作用。本文通过MuJoCo原生触觉传感器和后续处理获得712维标量观测，其评价需关注阵列的空间响应与幅值关系，与视觉触觉的成像评价有所区别。原生接触响应、阵列采样与处理后幅值应分别描述，接触求解中的柔顺性与真实敏感材料内部形变也属于不同的建模层次。对于同一接触，感知区域覆盖方式会影响哪些单元产生响应，滤波和邻域传播又会改变最终的激活范围与幅值。由此得到的空间扩散，既可能反映传感器模型，也可能来自读出后的处理，不能仅凭热图外观相近就认定材料响应一致。分析曲面触觉阵列时，需要保留几何布局、接触采样和信号处理之间的区别，才能判断差异出现在哪一环节。因而，传感器布局和处理过程的可追溯性，是分析仿真与真实触觉差异的前提。

\subsection{专家示教与多模态操作数据}
\label{subsec:demonstrations_multimodal_data}

操作示教可以通过人工遥操作获得，也可以依托规划、控制器和已有轨迹自动生成。人工能够根据执行反馈处理复杂情形，但扩大采集规模需要重复操作与场景复位；自动生成便于覆盖初态变化，却依赖任务描述、感知和控制接口。MimicGen将少量源示教分解为物体中心的子任务片段，经坐标变换和插值适配到新场景，再实际执行并保留成功轨迹\cite{mandlekar2023mimicgen}。其扩展能力建立在已知子任务顺序、物体位姿和末端控制形式等条件之上，不能直接推广到任意手型和操作。关于离线示教学习的系统性研究还表明，示教质量、观测选择及历史信息会影响策略表现\cite{mandlekar2022demonstrations}。因此，轨迹生成数量与轨迹对学习的实际价值需要分别检查。

数据集研究从采集来源、环境覆盖和模态组成等方面扩展了操作学习的条件。BridgeData V2组织跨环境、物体和任务的操作记录，并以目标及语言条件学习检验数据的用途\cite{walke2023bridgedata}；DROID采用分布式遥操作和统一采集装置拓展真实场景，并通过联合训练评价其利用效果\cite{khazatsky2024droid}。二者关注环境和操作多样性，上海交通大学的RH20T则进一步将视觉、力觉、音频和机器人状态等记录纳入采集与标定流程\cite{fang2024rh20t}。这些工作表明，不同数据集解决的采集问题并不相同。轨迹数和帧数描述数据规模，模态可用性、动作记录方式、校准和任务验证则规定了数据的适用范围。触觉观测的有无及其传感机制，仍需依据各数据集的实际组成分别判断。

跨机器人利用数据还涉及动作和观测语义的适配。Open X-Embodiment通过粗粒度的观测与动作整合支持跨机器人训练，但格式统一不意味着不同机器人的控制机制和采样条件完全一致\cite{openx2024embodiment}。对本文而言，同一仿真Episode内部的状态、动作和传感器记录，可以依据执行索引建立关联；来自不同真实设备的记录，还需要检查时间戳、坐标和处理流程。Real与Sim之间按统计描述构造的伪配对，只表示在所选统计量上的对应，而同步物理配对要求两域确有相应的接触过程。这些关联分别依赖执行索引、设备校准、统计匹配或物理过程，数组形状相同并不构成相互替代的依据。同一轨迹派生的相邻窗口通常存在相关性，重复导出也不会增加独立执行次数，因而数据规模还应结合采集单位和来源去重进行解释。示教生成需要与数据来源、时间关系和质量审计相配合，才能为后续重建研究提供含义明确的输入。

\subsection{触觉Sim-to-Real与跨域重建}
\label{subsec:tactile_sim_to_real}

Sim-to-Real研究需要处理模型近似和真实采集条件引起的差异。一类方法从仿真侧改善物理模型或扩大参数覆盖。Tobin等通过随机化纹理、光照和相机条件训练视觉定位模型，说明在特定任务中可以降低对单一渲染条件的依赖\cite{tobin2017randomization}；这一结果并不直接保证触觉材料或读出参数随机化同样有效。对触觉而言，传感器内部结构和敏感材料决定了接触到信号的转换。Sferrazza和D'Andrea利用有限元模型生成软材料形变和图像特征，以仿真力分布监督回归模型，并通过相机映射与校准应用于真实视觉触觉观测\cite{sferrazza2022simtoreal}。该路线以物理模型和标定条件连接两域，涉及的分布力真值来源也需与真实独立测量区分。因此，物理建模与数据驱动校准并非相互排斥，关键在于各阶段能够获得什么监督，以及模型误差会如何传递到最终读出。

与本文方法较为接近的是Narang等针对BioTac传感器开展的研究。他们建立三维有限元接触模型，分别以VAE学习仿真形变和真实电信号的潜表示，再利用有限的配对监督训练潜空间投影，实现两类观测之间的转换\cite{narang2021latentprojections}。其意义在于将复杂物理形变与传感器电响应分开表示，通过监督映射连接独立潜空间，而不是假定两套编码器自动对齐。本文面向途见科技灵巧手的712维处理后触觉，仿真端采用MuJoCo原生触觉读出，真实端使用既有实机记录，监督对应的构造条件与该研究不同。具体而言，利用Real Train拟合的GMM及Patch响应统计，在固定索引坐标下描述响应中心和协方差，再查询16个Patch的响应字典，确定球形探针的位置和加载条件，经MuJoCo实际接触生成与Real窗口统计对应的Paired Sim。这种对应属于统计伪配对，而非同步物理配对。Real VAE使用真实训练数据，Sim VAE同时使用Natural Sim和Paired Sim训练，两者独立学习表示；随后冻结两个VAE，利用Paired Sim与对应Real窗口的潜变量监督训练GRU Translator。推理时，Natural Sim依次经Sim Encoder、GRU Translator和Real Decoder生成完整四帧712维触觉，不再调用GMM或球形探针。因此，本文需要研究的是统计监督的可靠性及其向自然任务输入的适用性，而不是将双VAE与潜映射的组合本身视为新的研究贡献。

缺少配对监督时，迁移方法需要借助分布目标和额外结构约束。VAE通过重建与先验正则学习概率潜表示，但两域分别服从相同先验，并不保证潜变量各维具有相同语义\cite{kingma2014vae}。CycleGAN利用对抗目标和循环一致性限制非配对转换，UNIT则通过共享潜空间假设、权重共享及重建约束连接两域\cite{zhu2017cyclegan,liu2017unit}。这些方法为非配对重建提供了建模思路，同时也说明仅有两个边际分布不足以确定唯一对应。另一类域适应方法以任务表示为目标，例如DANN在保留源域任务判别性的同时，通过梯度反转抑制域信息\cite{ganin2016dann}。这种特征适应与完整触觉输出重建具有不同目的：前者关注任务所需表示，后者还需要恢复激活位置、幅值及帧间关系。潜空间中域信息的减少，为表示适应提供证据，完整712维观测是否对齐则仍需在输出端检验。

迁移评价也必须与上述目标对应。生成对抗网络关于分布一致的理论结论依赖理想判别器和优化条件，有限判别器下的高接受率并不自动满足这些前提\cite{goodfellow2014gan}。核两样本检验提供了比较分布的另一种方式，但其结果依赖核、表示和样本条件\cite{gretton2012mmd}；密度与覆盖等生成评价则分别检查样本质量和分布覆盖，仍受到特征空间选择的影响\cite{naeem2020metrics}。这些指标不能独自说明输出保留了哪一个输入的接触信息。训练中参与优化的判别器、迁移前冻结的域分类器和迁移后重新训练的分类器，检验的边界也不同，评价时应明确各自的数据与训练角色。对于触觉跨域重建，需要将判别器表现、分布统计和输入依赖分开考察，并检查空间响应与帧间变化。表\ref{tab:related_work_routes}归纳了相关路线的监督与适用条件。本文据此关注具体机器人上的可执行示教、明确来源的触觉读出，以及有限统计监督下的重建与结构保持，不预设某一表示或判别目标能够同时解决这些问题。

\begin{table}[htbp]
  \centering
  \caption{相关技术路线及其监督与适用条件}
  \label{tab:related_work_routes}
  \begingroup
  \setlength{\tabcolsep}{3pt}
  \begin{tabularx}{\textwidth}{p{0.10\textwidth}p{0.15\textwidth}XXX}
    \toprule
    技术方向 & 代表方法 & 主要输入与监督 & 解决的问题 & 适用条件及局限 \\
    \midrule
    抓取优化 & DexGraspNet、GraspQP & 物体和手模型、接触质量能量 & 多指抓取候选生成 & 依赖接触假设，不保证整机执行 \\
    物理细化 & DexEvolve & 解析候选、仿真反馈 & 稳定性与多样性优化 & 受物理模型及执行条件影响 \\
    目标执行 & UniDexGrasp & 点云、抓取目标和策略学习 & 连接候选与操作过程 & 依赖机器人与控制设置 \\
    触觉模拟 & TACTO、Taxim & 接触几何、光学或形变标定 & 视觉触觉读出生成 & 不能直接等同标量阵列响应 \\
    示教构建 & MimicGen及多模态数据集 & 源示教、场景和传感器记录 & 轨迹扩展与数据利用 & 需核查任务假设及数据关联 \\
    物理迁移 & BioTac潜投影\cite{narang2021latentprojections} & FEM形变、真实电信号和配对监督 & 跨模态潜表示转换 & 依赖监督对应和传感器模型 \\
    非配对迁移 & CycleGAN、UNIT & 两域样本及循环或共享约束 & 约束缺少配对的映射 & 不认证唯一物理对应 \\
    特征适应 & DANN & 源域任务标签及两域观测 & 任务相关的域不变表示 & 不等于完整触觉重建 \\
    \bottomrule
  \end{tabularx}
  \endgroup
\end{table}

\section{研究问题与技术难点}
\label{sec:research_problems}

前述研究表明，触觉示教数据的生成需要连接机器人操作与传感观测。基础问题是模型、控制与触觉阵列的一致性。驱动耦合约束手部运动，安装关系影响接近方向和接触位置；曲面阵列的单元位置、朝向与覆盖区域决定了接触信息如何被采样。原生接触读出经过死区、滤波等处理后，空间支持和幅值还可能发生变化。因此，需要明确各环节的输入输出与处理条件，使后续差异分析能够追溯到具体观测来源，而不是仅以712维形状一致认定物理含义相同。

在这一平台上，抓取生成需要解决候选配置与可执行过程之间的连接问题。几何目标或力闭合评价较好的候选，可能不满足手部驱动耦合，也可能使机械臂无法到达预抓取位置。即使两个目标位姿分别可达，连接它们的运动仍可能出现关节突变、碰撞或接触丢失。因而，候选筛选不仅涉及末端误差，还需要考虑接近、闭合、持握及搬运的连续过程，并按各任务要求检验终态与过程条件。需要连接各层约束并保留失败原因，分别界定抓取配置、任务成功与严格轨迹合格。

可执行轨迹只有与观测建立可靠关联，才能形成可用示教数据。状态、动作、视觉与触觉记录可能具有不同采样频率，观测发生在动作执行前还是执行后，也会改变其所表达的控制关系。转换中的索引丢失、重复与模态缺失，会使格式合法的数据仍存在内容问题。本文需要保留执行来源、时间和字段语义，将单次仿真内部关联与独立实机记录区分开，并通过一致性检查和质量审计识别异常。这里需要解决的是数据证据与操作过程之间的对应，而非仅完成文件保存或格式转换。

触觉跨域重建进一步受到监督条件限制。现有Real与Sim记录缺少经过认证的同步物理配对，真实统计特征能够约束响应比例、中心和扩散等描述，却不足以唯一确定产生这些响应的接触状态。由此构造的探针数据可能对应多个真实目标，其覆盖也可能不同于自然任务输入。两域独立学习的潜表示还存在坐标与语义差异，相同维度或先验不能替代对应监督。因此，需要区分伪配对构造、域内重建和潜变量映射分别承担的作用，考察有限统计监督能否支持完整触觉输出，而不将训练配对上的拟合直接推广到自然仿真数据。

重建结果还需要通过与目标相符的评价加以检验。某个域分类器接受生成样本，只能说明该判别边界下的读出发生变化；输出仍可能损失接触位置、幅值差异或帧间关系。分布接近也不自动证明生成结果依赖对应输入。真实采集会话之间的变化，还可能影响评价器对真实样本的接受与跨组泛化。本文因而需要分开检查域判别、分布质量和输入保持，并明确分类器的训练与冻结角色。评价器自身的分组表现和退化对照同样属于证据，据此判断优化改善所支持的结论范围。

\section{主要研究内容与创新点}
\label{sec:research_contents}

本文围绕仿真操作、示教数据和触觉重建开展研究。主要工作包括平台与观测集成、适应具体机器人约束的示教流程，以及有限统计监督下的跨域方法探索。三者分别提供实验条件、可追溯数据和重建分析依据，其贡献需要结合实现与验证范围说明，平台集成、开源适配和迁移效果不作同一层次的创新判断。上述三条主线对应四项工作：平台与触觉观测、抓取候选与专家示教、多模态示教数据与质量控制、有限监督触觉重建与评价。其中第二、三项同属专家示教与数据主线，第七章据此分项总结。

在平台与触觉建模方面，构建RM75B机械臂与途见科技自研欠驱动灵巧手的仿真组合，将任务环境、控制接口和曲面触觉阵列纳入操作流程。围绕感知单元的空间布局，建立原生触觉读取及后续信号处理链，获得712维标量观测，并在抓取抬升、抓取放置、堆叠、推动和螺母装配任务中形成开发验证记录。这部分工作的特色在于连接机器人运动、局部接触观测与任务记录，为不同操作阶段提供一致的仿真接口；其验证范围是模型与记录链路的功能，不等同于真实传感器物理标定或实机部署。

在专家示教与数据方面，结合GraspQP官方指标的欠驱动适配、DexEvolve-style物理搜索及腕部姿态扩展，构建多来源候选到机器人执行的分级流程。通过可达性筛选和轨迹回放衔接候选位姿与任务验证，并区分普通回放与严格检查的判定条件。在生成轨迹的基础上，组织状态、动作、视觉和触觉记录，开展格式转换、来源保留、重复识别及质量检查。工作重点是将第三方方法适配到具体驱动和执行约束，连接示教生成与数据审计，而非将原算法作为本文原创，也不以历史汇总代替全量数据验收。

在跨域重建方面，探索真实统计先验指导的仿真触觉伪配对方法，通过GMM统计描述和探针响应字典确定接触条件，经MuJoCo执行构造监督数据。分别学习真实域与仿真域VAE表示，在冻结两域模型后以GRU Translator学习潜变量映射，并经Real Decoder生成完整触觉。进一步研究共享参数向量场及迭代对抗优化，结合分布、空间、时序与输入依赖评价分析其效果。实验中，多轮优化提高了对应冻结判别器下的接受率，但质量未同步通过；多视角残差探索出现分布退化，后续审计则针对D50判定其不满足评价器使用前提。这些结果用于分析判别边界适应、分布逼近与结构保持之间的关系，不将局部接受率改善表述为全局迁移成功。

\section{论文组织结构}
\label{sec:thesis_structure}

全文按七章组织，由实验条件逐步过渡到数据构建、触觉重建及评价。第一章说明研究背景和相关技术，归纳研究问题及工作范围。第二章建立机械臂、欠驱动手和曲面触觉观测的仿真基础，交代控制与信号来源。第三章在此基础上研究抓取候选的适配、可达性筛选和轨迹验证，说明多任务专家示教如何形成。第四章从数据使用角度检查示教记录的组织、关联、转换和质量，使后续研究所使用的数据及其限制具有可追溯依据。第三、四章分别回答示教如何产生及如何成为可核验数据，避免将任务成功与数据质量混为一谈。

第五章围绕缺少同步物理配对的条件，介绍真实统计先验指导的探针伪配对、独立双域表示和潜变量映射，并分析初始重建的效果与局限。第六章承接初始迁移结果，研究共享参数向量场和迭代对抗优化，结合多视角判别、分布质量、输入保持及跨会话分析检验优化结果，同时讨论评价器的适用范围。第七章归纳各研究方向获得的结论，区分已实现功能、受限实验结果与尚未解决的问题，并提出与现有证据缺口相对应的后续方向。
